% !TITLE 旅夜书怀
% !AUTHOR 杜甫
% !DYNASTY 唐
% !ORDER 28

\begin{poem}
细草微风岸，危樯独夜舟\textsuperscript{1}。
星垂平野阔\textsuperscript{2}，月涌大江流\textsuperscript{3}。
名岂文章著，官应老病休\textsuperscript{4}。
飘飘何所似，天地一沙鸥\textsuperscript{5}。
\end{poem}

\begin{annotations}
\notetext{1}{危樯独夜舟：高高的桅杆，孤舟在夜中停泊。危，高；樯（qiáng），桅杆。}
\notetext{2}{星垂平野阔：星星低垂，原野显得格外辽阔。“垂”字写星星低挂欲坠，显得天地间更加空旷。}
\notetext{3}{月涌大江流：月光在江水中涌动，大江奔流不息。“涌”字写月影随波荡漾，有了动态。}
\notetext{4}{名岂文章著，官应老病休：名声哪里是因为文章显著的？官职也该因为年老多病而休退了。这是自嘲——杜甫不甘心只以文章出名，更希望建功立业。}
\notetext{5}{沙鸥：一只水鸟。天地之大，却孤零零的像一只沙鸥。这是杜甫漂泊一生的写照。}
\end{annotations}

\begin{appreciation}
这首诗写于唐代宗永泰元年（765年），杜甫离开成都草堂，沿江东下，在旅途中写下了这首感怀之作。

细草微风岸，危樯独夜舟——岸边长着细草，微风吹过；夜里一只孤舟，桅杆高耸。危樯就是高高的桅杆。独字是全诗的基调：一只船，一个人，一个夜晚。星垂平野阔，月涌大江流——星星低垂，原野显得格外辽阔；月亮在江水中涌动，大江滚滚东流。垂字写得好：星星不是在高空闪烁，而是低低地垂下来，像要落到地上一样。这种低垂让天空和大地之间的距离缩短了，也让孤独的人感到天地更加空旷。

名岂文章著，官应老病休——名声哪里是因为文章而显著的？官职也该因为年老多病而休退了。这是自嘲，也是不平。杜甫的文章确实好，但他不希望自己的名声仅仅来自文章——他更希望来自功业。官应老病休——不是因为不想做官了，是因为老了、病了，不得不休。一个应字，道出了多少无奈。

飘飘何所似，天地一沙鸥——漂泊无依像什么呢？就像天地间的一只沙鸥。这是全诗的结句，也是最孤独的一句。沙鸥是水鸟，孤零零地在天地间飞翔。它没有同伴，没有归宿，只有无边的天空和茫茫的水面。杜甫把自己比作沙鸥，不是一时的感慨，是一生的总结：他漂泊了一生，从洛阳到长安，从秦州到成都，从夔州到潭州——像一只永远找不到落脚之地的鸟。
\end{appreciation}
